\documentclass[10pt,a4paper]{amsart}
\usepackage[margin=19mm]{geometry}
\usepackage{amsmath,amssymb,mathtools}
\usepackage[scr=rsfs,cal=euler]{mathalfa}
\usepackage{xfrac}
\usepackage[unicode,hidelinks,bookmarks=false]{hyperref}
\newcommand{\ie}{i.e.\ }
\newcommand{\cf}{cf.\ }
\newcommand{\resp}{respectively\ }
\newcommand{\Subsection}{\subsection}
\providecommand{\nobreakdash}{\nobreak\hbox{-}}
\setlength{\emergencystretch}{2em}
\multlinegap0pt
\usepackage{bm}
\usepackage{bbm}
\usepackage{dsfont}
\usepackage{xspace}
\usepackage{cancel}
\usepackage{mathtools}
\usepackage{amsthm}
\makeatletter
\@ifundefined{lemma}{\newtheorem{lemma}{Lemma}}{}
\makeatother
\title[The Newton-Muon Optimizer]{The Newton-Muon Optimizer}
\author{Zhehang Du, Weijie Su}
\date{Source: arXiv:2604.01472v1 (CC BY 4.0)}
\begin{document}
\maketitle

\noindent\emph{Source and licence.} The Newton-Muon Optimizer. Version arXiv:2604.01472v1 (CC BY 4.0), distributed under the Creative Commons Attribution 4.0 International licence, as recorded in the arXiv licence metadata for this exact version and shown on the abstract page https://arxiv.org/abs/2604.01472v1. Full rights notice: licenses/arxiv-2604.01472-CC-BY-4.0.txt.\par\smallskip

\noindent\textit{Editorial scope.} Counted unit: the Lemma whose source label is \texttt{lem:ss-mode-wise} in the article (source0.tex lines 406--460, source label \texttt{lem:ss-mode-wise}), delivered together with the article's own complete original proof of it (source0.tex lines 462--525, one \texttt{proof} environment of 64 lines). No mathematics is written, repaired or re-derived: statement and proof are the article's own text line for line. Required context retained uncounted: eq:ss-grad (lines 346--349); eq:ss-Newton--Muon-update (lines 351--373); eq:ss-gd-update (lines 351--373); eq:ss-muon-update (lines 351--373); eq:ss-spike (lines 378--382); eq:ss-init-decomp (lines 384--391); eq:ss-residual-decomp (lines 395--402). Every definition, display and label of those lines is retained unchanged. Omitted from this package and not redistributed: 1--345, 350--350, 374--377, 383--383, 392--394, 403--405, 461--461, 526--1753. Nothing inside a retained excerpt is omitted or altered: every retained body is delivered exactly as the article's lines are written, so no omission receipt of any kind accompanies this package. Citations: the retained excerpts carry no citation command at all, so no bibliography entry of the article is transcribed and no reference is redistributed. Reference adaptation: none was needed and none was made; every reference inside the delivered unit resolves inside it, so no unresolved reference or citation is produced. Printed slips kept literal: no printed word, symbol, hypothesis or number of the counted statement or of its proof was altered. Layout and class: the article's own class is \texttt{article} with packages including amsmath, amssymb, amsthm, mathrsfs, mathtools, algorithm2e, geometry, float, natbib, graphicx, appendix, tikz-cd, xcolor, xurl; the wrapper is the lane's pinned \texttt{amsart} editorial wrapper at 10pt on A4, which restates the theorem-like environments the retained text needs and adds the article's own macro definitions that the retained text needs (none), with extra wrapper packages (bm, bbm, dsfont, xspace, cancel, mathtools). The delivered numbering is the unit's own. Assets: no retained span contains a figure environment, a graphics-inclusion command, a PDF-page inclusion, a table or a TikZ picture; no image file is copied into this package, the package declares no asset dependency and no image file is redistributed with it. Licence: version arXiv:2604.01472v1 of this article is distributed under the Creative Commons Attribution 4.0 International licence, as recorded in the arXiv licence metadata for this exact version and shown on the abstract page https://arxiv.org/abs/2604.01472v1. A full rights notice is delivered at licenses/arxiv-2604.01472-CC-BY-4.0.txt. Characters: every retained excerpt is pure ASCII, so the delivered engine, XeLaTeX with its default Latin Modern text font, reports no missing character.
\begin{equation}
  \label{eq:ss-grad}
  \nabla_{{W}} f({W}) = ({W}-W^\star){Z}{Z}^\top.
\end{equation}

\begin{align}
  \label{eq:ss-gd-update}
  \text{GD:}\quad
  W_{t+1}
  &=
  W_t-\eta_t \nabla_{W} f(W_t)
  =
  W_t-\eta_t ( W_t-W^\star )ZZ^\top, \\
  \label{eq:ss-muon-update}
  \text{Muon:}\quad
  W_{t+1}
  &=
  W_t-\eta_t \mathrm{msgn}\big(\nabla_{W} f(W_t)\big)
  =
  W_t-\eta_t \mathrm{msgn}\big(( W_t-W^\star )ZZ^\top\big), \\
  \label{eq:ss-Newton--Muon-update}
  \text{Newton--Muon:}\quad
  W_{t+1}
  &=
  W_t-\eta_t \mathrm{msgn}\big(\nabla_{W} f(W_t)(ZZ^\top)^{-1}\big)
  =
  W_t-\eta_t \mathrm{msgn}(W_t-W^\star).
\end{align}

\begin{equation}
  \label{eq:ss-spike}
  {Z}{Z}^\top = {U}_{Z}\mathrm{diag}(\kappa,1,\ldots,1){U}_{Z}^\top,
  \qquad \kappa>1,
\end{equation}

\begin{equation}
  \label{eq:ss-init-decomp}
  W^\star
  =
  -\boldsymbol{u}_1 \big(\alpha_{1,0} \boldsymbol{e}_1^\top + \beta_{1,0} \boldsymbol{b}_1^\top\big)
  -
  \sum_{i=2}^r \boldsymbol{u}_i \big(\beta_{i,0} \boldsymbol{b}_i^\top\big),
\end{equation}

\begin{equation}
  \label{eq:ss-residual-decomp}
  W_t-W^\star
  =
  \boldsymbol{u}_1 \big(\alpha_{1,t} \boldsymbol{e}_1^\top + \beta_{1,t} \boldsymbol{b}_1^\top\big)
  +
  \sum_{i=2}^r \boldsymbol{u}_i \big(\beta_{i,t} \boldsymbol{b}_i^\top\big).
\end{equation}

\begin{lemma}[Dynamics under the single spike model]
\label{lem:ss-mode-wise}
Under \eqref{eq:ss-spike} and \eqref{eq:ss-init-decomp}, gradient descent \eqref{eq:ss-gd-update}, Muon \eqref{eq:ss-muon-update}, and Newton--Muon \eqref{eq:ss-Newton--Muon-update} preserve the decomposition \eqref{eq:ss-residual-decomp}. In particular, the dynamics decouple across modes $i=1,\dots,r$, and the coefficients satisfy the following recursions.
\begin{itemize}
\item For gradient descent,
\begin{equation}
\label{eq:ss-rec-gd}
\alpha_{1,t+1}=(1-\eta_t\kappa)\alpha_{1,t},
\qquad
\beta_{i,t+1}=(1-\eta_t)\beta_{i,t},\quad i=1,\dots,r.
\end{equation}

\item For Muon,
\begin{equation}
\label{eq:ss-rec-muon}
\alpha_{1,t+1}
=
\alpha_{1,t}
-
\eta_t \frac{\kappa\alpha_{1,t}}{\sqrt{\kappa^2\alpha_{1,t}^2+\beta_{1,t}^2}},
\qquad
\beta_{i,t+1}
=
\begin{cases}
\beta_{1,t}
-\eta_t \bigl(\beta_{1,t}/\sqrt{\kappa^2\alpha_{1,t}^2+\beta_{1,t}^2}\bigr),
& i=1, \\
\beta_{i,t}
-\eta_t \bigl(\beta_{i,t}/|\beta_{i,t}|\bigr),
& i=2,\dots,r.
\end{cases}
\end{equation}

\item For Newton--Muon,
\begin{equation}
\label{eq:ss-rec-Newton--Muon}
\alpha_{1,t+1}
=
\alpha_{1,t}
-
\eta_t \frac{\alpha_{1,t}}{\sqrt{\alpha_{1,t}^2+\beta_{1,t}^2}},
\qquad
\beta_{i,t+1}
=
\begin{cases}
\beta_{1,t}
-\eta_t \bigl(\beta_{1,t}/\sqrt{\alpha_{1,t}^2+\beta_{1,t}^2}\bigr),
& i=1, \\
\beta_{i,t}
-\eta_t \bigl(\beta_{i,t}/|\beta_{i,t}|\bigr),
& i=2,\dots,r.
\end{cases}
\end{equation}
\end{itemize}
\end{lemma}

\begin{proof}[Proof of Lemma~\ref{lem:ss-mode-wise}]
Define $\boldsymbol{v}_{1,t}\coloneqq \alpha_{1,t}\boldsymbol{e}_1+\beta_{1,t}\boldsymbol{b}_1$ and, for $i\ge 2$, $\boldsymbol{v}_{i,t}\coloneqq \beta_{i,t}\boldsymbol{b}_i$, so that
\begin{equation}
  \label{eq:ss-v-decomp}
  W_t-W^\star=\sum_{i=1}^r \boldsymbol{u}_i\boldsymbol{v}_{i,t}^\top.
\end{equation}
Since $\{\boldsymbol{u}_i\}_{i=1}^r$ are orthonormal and $\{\boldsymbol{b}_i\}_{i=1}^r$ are orthonormal with $\boldsymbol{b}_i\perp \boldsymbol{e}_1$, we have $\boldsymbol{v}_{i,t}\perp \boldsymbol{v}_{j,t}$ for $i\neq j$ for every $t$.
Under \eqref{eq:ss-spike}, $ZZ^\top \boldsymbol{e}_1=\kappa \boldsymbol{e}_1$, and the eigenspace with eigenvalue $1$ is $\boldsymbol{e}_1^\perp$, so $ZZ^\top \boldsymbol{b}_i=\boldsymbol{b}_i$ for all $i=1,\dots,r$. Hence
\begin{equation}
  \label{eq:ss-v-mapped}
  ZZ^\top \boldsymbol{v}_{1,t}=\kappa\alpha_{1,t}\boldsymbol{e}_1+\beta_{1,t}\boldsymbol{b}_1 \in \mathrm{span}\{\boldsymbol{e}_1,\boldsymbol{b}_1\},
  \qquad
  ZZ^\top \boldsymbol{v}_{i,t}=\boldsymbol{v}_{i,t}\in \mathrm{span}\{\boldsymbol{b}_i\}\ \ (i\ge 2).
\end{equation}
Moreover, the transformed vectors $\{ZZ^\top\boldsymbol{v}_{i,t}\}_{i=1}^r$ remain mutually orthogonal across $i$. Using \eqref{eq:ss-grad} and \eqref{eq:ss-v-decomp},
\begin{equation}
  \label{eq:ss-grad-decomp}
  \nabla_W f(W_t)=(W_t-W^\star)ZZ^\top
  =
  \sum_{i=1}^r \boldsymbol{u}_i\big((ZZ^\top\boldsymbol{v}_{i,t})^\top\big).
\end{equation}
In particular, by \eqref{eq:ss-v-mapped},
\[
(ZZ^\top\boldsymbol{v}_{1,t})^\top=(\kappa\alpha_{1,t})\boldsymbol{e}_1^\top+\beta_{1,t}\boldsymbol{b}_1^\top,
\qquad
(ZZ^\top\boldsymbol{v}_{i,t})^\top=\beta_{i,t}\boldsymbol{b}_i^\top \ (i\ge 2).
\]

For gradient descent, substituting \eqref{eq:ss-grad-decomp} into \eqref{eq:ss-gd-update} gives
\[
W_{t+1}-W^\star
=
(W_t-W^\star)-\eta_t(W_t-W^\star)ZZ^\top,
\]
and matching coefficients in the basis $\{\boldsymbol{u}_1\boldsymbol{e}_1^\top,\boldsymbol{u}_1\boldsymbol{b}_1^\top,\boldsymbol{u}_i\boldsymbol{b}_i^\top\}$ yields $\alpha_{1,t+1}=(1-\eta_t\kappa)\alpha_{1,t}$ and $\beta_{i,t+1}=(1-\eta_t)\beta_{i,t}$ for all $i=1,\dots,r$, i.e., \eqref{eq:ss-rec-gd}.

For Muon, since the right factors $\{ZZ^\top\boldsymbol{v}_{i,t}\}_{i=1}^r$ in \eqref{eq:ss-grad-decomp} are mutually orthogonal and the left factors $\{\boldsymbol{u}_i\}_{i=1}^r$ are orthonormal, the matrix sign of $(W_t-W^\star)ZZ^\top$ is given by
\[
\mathrm{msgn}\big((W_t-W^\star)ZZ^\top\big)
=
\sum_{i:\|ZZ^\top\boldsymbol{v}_{i,t}\|_2>0}
\boldsymbol{u}_i
\left(\frac{(ZZ^\top\boldsymbol{v}_{i,t})^\top}{\|ZZ^\top\boldsymbol{v}_{i,t}\|_2}\right).
\]
For $i\ge 2$, $ZZ^\top\boldsymbol{v}_{i,t}=\boldsymbol{v}_{i,t}$ and $\|ZZ^\top\boldsymbol{v}_{i,t}\|_2=|\beta_{i,t}|$, yielding the $i\ge 2$ case in \eqref{eq:ss-rec-muon}. For $i=1$, we have
$ZZ^\top\boldsymbol{v}_{1,t}=\kappa\alpha_{1,t}\boldsymbol{e}_1+\beta_{1,t}\boldsymbol{b}_1$ and
$\|ZZ^\top\boldsymbol{v}_{1,t}\|_2=\sqrt{\kappa^2\alpha_{1,t}^2+\beta_{1,t}^2}$, so substitution into \eqref{eq:ss-muon-update} yields the $i=1$ case in \eqref{eq:ss-rec-muon}. This also preserves \eqref{eq:ss-residual-decomp}.

For Newton--Muon, since the right factors $\{\boldsymbol{v}_{i,t}\}_{i=1}^r$ are mutually orthogonal and the left factors $\{\boldsymbol{u}_i\}_{i=1}^r$ are orthonormal, the matrix sign of $W_t-W^\star$ is obtained mode-by-mode:
\[
\mathrm{msgn}(W_t-W^\star)
=
\sum_{i:\|\boldsymbol{v}_{i,t}\|_2>0}
\boldsymbol{u}_i
\left(\frac{\boldsymbol{v}_{i,t}}{\|\boldsymbol{v}_{i,t}\|_2}\right)^\top.
\]
Substituting into \eqref{eq:ss-Newton--Muon-update} gives
\[
W_{t+1}-W^\star
=
(W_t-W^\star)-\eta_t\mathrm{msgn}(W_t-W^\star),
\]
and matching coefficients yields the $i=1$ and $i\ge 2$ cases in \eqref{eq:ss-rec-Newton--Muon}. This preserves \eqref{eq:ss-residual-decomp} and completes the proof.
\end{proof}

\end{document}
